% BHU PYQ -- Transcribed & Corrected
% Paper: CHEMJ41: Inorganic Chemistry-II
% Exam: B.Sc. (Hons) Semester IV Examination 2025-26 (NEP)
% Ref Code: Conf/Even/N/2025-26/05
% Category: Major
% Department: Chemistry

\documentclass[11pt,a4paper]{article}
\usepackage[a4paper,top=2.5cm,bottom=2.5cm,left=2.8cm,right=2.8cm,headheight=14pt]{geometry}
\usepackage{amsmath,amssymb}
\usepackage[T1]{fontenc}
\usepackage[utf8]{inputenc}
\usepackage{lmodern,microtype,fancyhdr,enumitem,hyperref,lastpage,parskip}

\hypersetup{
    colorlinks=true,
    urlcolor=black,
    linkcolor=black,
    pdftitle={CHEMJ41: Inorganic Chemistry-II -- BHU Sem IV 2025-26 (NEP)}
}

\pagestyle{fancy}
\fancyhf{}
\renewcommand{\headrulewidth}{0.4pt}
\renewcommand{\footrulewidth}{0.4pt}
\fancyhead[L]{\small \textit{Banaras Hindu University}}
\fancyhead[R]{\small \textit{CHEMJ41: Inorganic Chemistry-II (2025-26)}}
\fancyfoot[C]{\small Page \thepage\ of \pageref{LastPage}}

\newcommand{\pts}[1]{\hfill \textit{[#1]}}

\begin{document}

\begin{center}
  {\large\bfseries BANARAS HINDU UNIVERSITY}\\[2pt]
  {\normalsize B.Sc. (Hons) Semester IV Examination 2025-26 (NEP)}\\[6pt]
  \rule{0.8\linewidth}{0.5pt}\\[6pt]
  {\large\bfseries CHEMISTRY}\\[4pt]
  {\large\bfseries Paper Code: CHEMJ41 : Inorganic Chemistry-II}\\[2pt]
  {\small Ref. No.: Conf/Even/N/2025-26/05}\\[4pt]
  {\normalsize Time: 2$\frac{1}{2}$ Hours\hfill Maximum Marks: 70}
\end{center}

\rule{\linewidth}{0.4pt}

\smallskip

\noindent \textit{Instructions: Answer five questions including Question 1 which is compulsory. The figures in the right-hand margin indicate marks.}

\bigskip

\noindent \textbf{Question 1.}
\begin{enumerate}[label=(\alph*)]
  \item CsI is much less soluble in water than CsF and LiF is much less soluble than LiI. Why? \pts{2}
  \item Why is $[\text{Ni(en)}_3]^{2+}$ more stable than $[\text{Ni(NH}_3)_6]^{2+}$ though both are octahedral complexes? \pts{2}
  \item What is the difference between solvation and solvolysis? Explain with one example of each type. \pts{2}
  \item What is E-factor in green chemical processes? \pts{2}
  \item What is lanthanide contraction and what are its important consequences? \pts{2}
  \item Write the names and electronic configurations of two naturally occurring actinide elements. \pts{2}
  \item What is Hammett acidity function? \pts{2}
\end{enumerate}

\bigskip

\noindent \textbf{Question 2.}
\begin{enumerate}[label=(\alph*)]
  \item Draw the various stereoisomers and specify the enantiomeric pairs in (i) $\text{Ma}_4\text{b}_2$ and (ii) $\text{M(AA)}_3$ type of octahedral complexes. \pts{6}
  \item State the properties that determine the utility of a solvent. \pts{4}
  \item What are the differences in characteristic properties of blue colored solution and bronze colored phase obtained by addition of alkali metals in liquid ammonia? \pts{4}
\end{enumerate}

\bigskip

\noindent \textbf{Question 3.}
\begin{enumerate}[label=(\alph*)]
  \item Briefly write about the various principles of green chemistry. \pts{6}
  \item Arrange the following acids in decreasing order of acid strength in aqueous solution. Explain the reason behind it: $\text{H}_2\text{SO}_3$, $\text{H}_2\text{SO}_4$, $\text{HMnO}_4$, $\text{H}_3\text{AsO}_4$. \pts{2}
  \item Give two examples each of polar aprotic and polar protic solvents and write autoionization reaction for them. \pts{3}
  \item Predict the product when $(\text{CH}_3)_2\text{N-PF}_2$ is reacted with (i) $\text{B}_2\text{H}_6$ and (ii) $\text{BF}_3$. Explain the reactivity difference in the two cases. \pts{3}
\end{enumerate}

\bigskip

\noindent \textbf{Question 4.}
\begin{enumerate}[label=(\alph*)]
  \item Using valence bond theory, predict the structure and spin only magnetic moment of (i) potassium tetracyanidonickelate(II) and (ii) hexaamminecobalt(III) chloride. \pts{4}
  \item Write the electronic configuration and number of unpaired electrons in the ground state of following ions: $\text{Tb}^{4+}$, $\text{Gd}^{3+}$, $\text{Ce}^{4+}$, $\text{Lu}^{3+}$. \pts{4}
  \item Explain the reason behind following observations: \pts{4}
  \begin{enumerate}[label=(\roman*)]
    \item Silver chloride is colorless while silver iodide is yellow in color.
    \item $[\text{CoCl}_4]^{2-}$ is more intensely colored than $[\text{MnBr}_4]^{2-}$.
  \end{enumerate}
  \item Complete the following reactions and identify acid and base in each case: \pts{2}
  \begin{enumerate}[label=(\roman*)]
    \item $\text{NH}_2\text{SO}_2\text{OH} + \text{H}_2\text{O} \rightarrow ?$
    \item $\text{CaO} + \text{SiO}_2 \rightarrow ?$
  \end{enumerate}
\end{enumerate}

\bigskip

\noindent \textbf{Question 5.}
\begin{enumerate}[label=(\alph*)]
  \item Identify and draw the various structural isomers of the following complexes: \pts{6}
  \begin{enumerate}[label=(\roman*)]
    \item $[\text{Cr(NH}_3)_6][\text{Co(CN)}_6]$
    \item $[\text{CoBr(NH}_3)_5]\text{SO}_4$
    \item $[\text{Cr(H}_2\text{O})_6]\text{Cl}_3$
  \end{enumerate}
  \item What are the limitations of effective atomic number rule? \pts{2}
  \item What are the various methods of separating the lanthanides? \pts{6}
\end{enumerate}

\bigskip

\noindent \textbf{Question 6.}
\begin{enumerate}[label=(\alph*)]
  \item What are intramolecular and intermolecular frustrated Lewis acid-base pairs? Explain with an example of each type. \pts{4}
  \item Explain the following with the help of suitable examples: \pts{8}
  \begin{enumerate}[label=(\roman*)]
    \item Advantage of solvent free reaction in green chemical synthesis.
    \item Difference between ordinary acids and superacids.
  \end{enumerate}
  \item The most common source of mercury is cinnabar ($\text{HgS}$), whereas $\text{Zn}$ and $\text{Cd}$ in the same group occur more commonly as carbonate, silicate and oxide. Why? \pts{2}
\end{enumerate}

\bigskip

\noindent \textbf{Question 7.}
\begin{enumerate}[label=(\alph*)]
  \item Briefly explain the various methods of uranium enrichment. \pts{6}
  \item What are ionic liquids? Briefly classify them. \pts{4}
  \item Write the IUPAC names of the following: \pts{4}
  \begin{enumerate}[label=(\roman*)]
    \item $\text{K}_3[\text{Fe(CN)}_6]$
    \item $\text{K}[\text{Cr(C}_2\text{O}_4)_2(\text{H}_2\text{O})_2]$
  \end{enumerate}
  \item Draw the structure of: $[\text{Co(NO}_2)(\text{NH}_3)_5]\text{Cl}_2$ and $\text{Fe}_2(\text{CO})_9$.
\end{enumerate}

\bigskip
\begin{center}
  \textbf{***}
\end{center}

\end{document}
